\subsection{MPRA data and fixed chromosome partitions}

We used the public MPRA sequence-performance table from Gosai et
al.~\citep{Gosai2024}, containing activity estimates and standard errors for K562,
HepG2, and SK-N-SH. The upstream chromosome partition was preserved: chromosomes 19,
21, and X were validation; chromosomes 7 and 13 were test; remaining eligible
chromosomes were training. Records were required to contain only A/C/G/T, have length
at most 200 nt, contain finite activity and standard-error values, and have maximum
reported log-fold-change standard error below 1.0. Exact sequence duplicates were
removed before full-data model training. Sequences shorter than 200 nt were centered
in a 200-position tensor, and a fifth input channel marked valid positions.

\subsection{Leakage and parent-overlap auditing}

We examined exact identity, reverse-complement identity, primary-variant group
membership, and cross-split Hamming distance up to three among 200-nt constructs. The
near-duplicate audit identified 27 cross-split pairs involving eight training-side
variant groups. Parent deny lists included earlier development and confirmation
parents, matching identifiers, sequence hashes, and audit-flagged conflicts. The
600-parent primary-evaluation table was created without activity labels by sorting
eligible 200-nt test sequences by SHA-256 of \texttt{ID + tab + sequence} and
selecting the first 600 after deny-list removal. The ablation study used a fixed
180-parent subset, and the nine-criterion benchmark used a fixed 96-parent subset.
Parent overlap across phases was checked by identifier and sequence hash.

The full-data loader performed exact deduplication. The separate Hamming-distance
manifest was applied to parent eligibility, ensuring that no flagged near-duplicate
was admitted to the editing benchmarks.

\subsection{Primary predictor}

Malinois BassetBranched, released with Gosai et al.~\citep{Gosai2024}, served as the
primary predictor. The checkpoint was loaded in weights-only mode. Each sequence was
evaluated in forward and reverse-complement orientations and the activity prediction
was averaged. For target cell $t$, the specificity margin was
\begin{equation}
M_t(x)=f_t(x)-\max_{j\ne t} f_j(x),
\end{equation}
where $f_j(x)$ is predicted activity in cell type $j$. Candidate margin gain was
$\Delta M_t=M_t(x_{\mathrm{edit}})-M_t(x_{\mathrm{parent}})$.

\subsection{Reviewer and cross-model evaluation ensembles}

Two architecture-diverse heteroscedastic CNNs were trained on all quality-filtered
training records to provide independent evaluation at different stages. The reviewer
ensemble is a residual-dilated 1D CNN with an 11-nt stem, residual blocks at
dilations 1, 2, 4, and 8, and shared mean and log-variance heads; it is available
during candidate search for cross-model reranking. The cross-model evaluation
ensemble (termed the post-freeze evaluator in the procedural protocol below) is a
multi-kernel CNN with parallel 3-, 7-, 15-, and 21-nt branches, pointwise fusion,
depthwise convolution, and separate mean and variance
heads~\citep{He2016,Lakshminarayanan2017}. Each architecture used three fixed seeds:
20260721--20260723 for the reviewer and 20260731--20260733 for the cross-model
evaluation ensemble.

Models were trained for at most 20 epochs with AdamW~\citep{Loshchilov2019}, cosine
learning-rate decay, batch size 256, learning rate $10^{-3}$, weight decay $10^{-4}$,
dropout 0.25, and gradient clipping at 5.0. The loss combined heteroscedastic Gaussian
negative log likelihood, reported measurement standard error, and a
reverse-complement consistency penalty weighted 0.1. Early stopping used mean
validation Pearson correlation with patience six. Test records were not used to
choose checkpoints, architectures, hyperparameters, or variance scales.

For each ensemble, prediction was the mean of the three seed-specific means, in the
spirit of bootstrap-aggregated ensemble prediction \citep{Breiman2001}. Predictive
variance combined mean learned within-model variance with between-seed variance. One
multiplicative variance scale per cell type was estimated from squared errors on
validation only and then frozen~\citep{Ovadia2019}. The cross-model evaluation
ensemble was never imported by the search code; it was applied only after candidate
files and their SHA-256 hashes had been written (procedurally frozen), ensuring
complete separation between generation and evaluation.

\subsection{Editing methods}

All methods were evaluated for target cells K562, HepG2, and SK-N-SH at exact
substitution budgets 1, 5, 10, and 20. Positions could be edited at most once along a
search path.

\paragraph{Random matched edits.}
Positions and non-reference alleles were sampled uniformly without replacement.
The 600-parent evaluation used three independent fixed seeds (20260801--20260803).
Numeric endpoints were averaged across replicates for method-level comparisons; the
best random replicate was never selected.

\paragraph{Greedy Malinois.}
At each step, every available single-nucleotide substitution was scored by the primary
specificity margin and the highest-scoring substitution was retained. No review,
uncertainty, or sequence-domain constraints were enforced during search.

\paragraph{Primary beam (compute-matched control).}
A beam of 24 partial edit sets was expanded one substitution at a time. Up to 96
candidates per step were retained by the primary score, which was the primary margin
minus 0.05 times forward/reverse strand disagreement. Beam width, maximum expansion
depth, and pre-screen size exactly matched SafeEdit-CRE. No reviewer score,
uncertainty penalty, or sequence-domain reranking was used.

\paragraph{SafeEdit-CRE.}
SafeEdit-CRE used the same primary expansion and compute budget as the primary beam.
The top primary candidates were rescored by the reviewer ensemble according to
\begin{equation}
S(x)=M_t^{\mathrm{primary}}(x)+0.3M_t^{\mathrm{reviewer}}(x)
-0.1D_{\mathrm{strand}}(x)-0.05U_{\mathrm{reviewer}}(x).
\end{equation}
Candidates were prefiltered for homopolymer length above 8 and absolute parent-relative
GC change above 0.12. Reranking used validation-derived, budget-specific limits for
strand disagreement, reviewer uncertainty, 6-mer naturalness change, GC change, and
homopolymer length. Positive primary and reviewer margins were favored; fallback
scores retained a path when no candidate passed every preferred condition. A snapshot
was taken independently at each requested budget.

Sequence naturalness was the position-independent mean log likelihood of overlapping
6-mers under frequencies estimated from natural training sequences. It was used as a
computational sequence-domain diagnostic during reranking.

\subsection{Nine-criterion design benchmark}

The nine-criterion benchmark (G4) used 96 test-split CRE parents with no overlap to
development or primary-evaluation parents. For each method--parent--target--budget
combination, nine Boolean checks were recorded: non-negative primary target gain,
positive primary margin gain, positive reviewer margin transfer, strand disagreement
within domain, reviewer uncertainty within domain, naturalness within domain, GC
change within domain, homopolymer safety, and exact edit budget. Thresholds were
calibrated from validation-set random edits and locked before evaluation. Infeasible
combinations were required to remain in the result table with failed status.

The primary contrast was the paired nine-criterion pass-rate difference between
SafeEdit-CRE and greedy. We used a 100,000-replicate parent-cluster bootstrap for 95\%
confidence intervals, preserving the dependence among the 12 target--budget
conditions contributed by each parent. An exact two-sided McNemar test on pooled
binary outcomes was retained as a descriptive sensitivity analysis. Stratified
target--budget summaries were used to examine heterogeneity.

\subsection{Primary cross-model evaluation (600 parents)}

The primary evaluation (G8) evaluated 600 label-blind, previously unused parents with
random matched, greedy, primary beam, and SafeEdit-CRE methods. Search produced 43,200
sequence rows because random editing had three replicates. For method-level
statistics, the three random replicates were averaged to one numeric record per
parent--target--budget combination, giving 28,800 rows. This collapsed table was used
only for numeric comparisons, never to select a DNA sequence. Candidate files were
frozen (SHA-256 hashes written) before cross-model-ensemble inference.

The principal endpoint was cross-model specificity-margin gain. Secondary endpoints
were cross-model target gain, positive-margin fraction, and cross-model ensemble
uncertainty. For each contrast, metrics were first averaged over the 12
target--budget conditions within a parent; 100,000 bootstrap samples then resampled
the 600 parents with replacement. Infeasible beam-search outcomes were retained as
parent sequences with zero edits and zero gain.

\subsection{Complete re-search ablation}

Seven searches were independently rerun on 180 parents: primary beam, full
SafeEdit-CRE, SafeEdit-CRE without reviewer score, without uncertainty penalty,
without strand penalty, without naturalness reranking, and without the sequence
prefilter. Unlike post-hoc filter deletion, each ablation could generate a different
sequence. The cross-model evaluation ensemble was then rerun on all 15,120 candidate
rows. We verified candidate keys, exact edit budgets, finite numeric values, and
separate computation of cross-model target and margin gains. Paired parent-cluster
bootstrap intervals used the same procedure as the primary evaluation.

\subsection{Candidate tiering}

The 3,456 nine-criterion-benchmark candidates were organized into three computational
priority tiers. Tier A required SafeEdit-CRE generation, passage of all nine
criteria, non-negative target gain, positive primary and reviewer margin gains, exact
edit budget, no synthesis-risk flags, and location on or near the two-objective
primary--reviewer Pareto frontier. Near-Pareto candidates were within 5\% relative or
0.1 absolute of the frontier. Tier B passed core activity and cross-model checks but
permitted Pareto suboptimality or mild naturalness deviation. Tier C failed one or
more hard checks and is retained only for diagnostic analysis. Three representative
Tier A examples spanning K562, HepG2, and SK-N-SH are provided in Supplementary
Table~S4.

\subsection{Software, reproducibility, and reporting}

Preliminary development phases were run on an NVIDIA RTX 3060 Laptop GPU. Final
benchmarks were run on RTX 3080 Ti compute nodes under Slurm. Fixed seeds, environment
records, code hashes, candidate hashes, checkpoint hashes, and run logs are provided
in the permanent Zenodo archive (v1.0). Statistical values and figures were
regenerated directly from candidate-level tables by an independent analysis script
implemented with PyTorch, NumPy, pandas, and
Matplotlib~\citep{Paszke2019,Harris2020,McKinney2010,Hunter2007}. All code,
environment files, minimal reproduction scripts, core statistical inputs, candidate
sequence tables, and pretrained model checkpoints are included in the public release.
